\documentclass[11pt]{article}
\usepackage{lmodern}
\usepackage[T1]{fontenc}
\usepackage[margin=1in]{geometry}
\usepackage{longtable}
\usepackage{array}
\usepackage{booktabs}
\usepackage{xcolor}
\usepackage{hyperref}
\renewcommand{\arraystretch}{1.3}
\setlength{\LTleft}{0pt}\setlength{\LTright}{0pt}

\begin{document}

\begin{center}
{\large\bfseries Response to Reviewers' Comments}\\[2pt]
Paper ID 10: \emph{Hydration Leakage and Over-fetching in Next.js:
The Security Implications of Server-to-Client Data Serialization}\\[2pt]
Arnav Garg, BITS Pilani (CAIS'26)
\end{center}

\vspace{4pt}
\noindent We thank the reviewers and editors for their careful feedback.
Below, each comment is paired with the corresponding revision.

\vspace{6pt}
\begin{longtable}{@{}p{0.44\textwidth}p{0.52\textwidth}@{}}
\toprule
\textbf{Reviewer / editor comment} & \textbf{Author's response and action taken} \\
\midrule
\endfirsthead
\multicolumn{2}{@{}l}{\emph{(continued)}}\\
\toprule
\textbf{Reviewer / editor comment} & \textbf{Author's response and action taken} \\
\midrule
\endhead
\bottomrule
\endfoot

\multicolumn{2}{@{}l}{\textbf{Reviewer \#1}}\\
\midrule
1. All abbreviations should be predefined. &
Done. Every abbreviation (SSR, HTML, API, SSN, RSC, PoC, DTO, JSON, HTTP,
HTTPS, CSP, SQL, TOTP, JWT, PAT, AWS, ORM, CWE) is now expanded on first
use. \\

2. A comparison with related works could be presented. &
Added \S2.5 ``Comparison with Existing Tools'' and Table~1, comparing the
scanner with the React taint API, secret scanners, generic SAST, and
architectural excessive-exposure metrics. \\

3. False positive rate is not quantified. &
Added \S4.3 ``Detection Effectiveness'' with a full confusion matrix
(Table~6): precision, recall, and $F_1$. We report recall 100\%, raw
precision 20.5\%, and 100\% after a short internal-field allowlist. \\

4. No evaluation against real-world Next.js applications. &
Added a cross-framework evaluation (\S4.4, Table~7) and a
``Threats to Validity'' subsection (\S6.4) that scopes an authorized,
at-scale study of deployed applications as future work, with the
legal/ethical rationale. \\
\midrule

\multicolumn{2}{@{}l}{\textbf{Reviewer \#3}}\\
\midrule
Additional experiments on larger/real-world datasets. &
Added the cross-framework experiment and a scalability benchmark
(\S4.5, up to 10{,}000-record payloads); production-scale scanning is
scoped as future work in \S6.4. \\

Add architectural / workflow diagrams. &
Added Fig.~1 (hydration serialization pipeline) and Fig.~2 (scanner
workflow flowchart). \\

Discuss scanner performance (runtime, scalability, detection efficiency). &
Added \S4.5 ``Performance and Scalability'': linear $O(n)$ time,
$\approx$2.1\,$\mu$s per node, a 2.5\,MB payload analyzed in $\approx$200\,ms. \\

Comparison with existing security/static-analysis tools. &
Added \S2.5 and Table~1 (see Reviewer \#1, comment 2). \\

Expand related work with academic references. &
Added five peer-reviewed references and rewrote \S2.2--2.4 around them. \\
\midrule

\multicolumn{2}{@{}l}{\textbf{Reviewer \#9}}\\
\midrule
Additional evaluation on larger real-world apps / public datasets. &
Addressed via the cross-framework experiment (\S4.4) and the scalability
benchmark (\S4.5); at-scale production study scoped as future work. \\

Compare the scanner with existing tools. &
Added \S2.5 and Table~1. \\

Add architectural diagrams. &
Added Figs.~1 and 2. \\

Expand applicability to other SSR frameworks (Nuxt, Remix, SvelteKit). &
The scanner was extended to all three; \S4.4 (Table~7) reports 32/32
credentials detected across the four frameworks, and \S6.2 confirms the
DTO mitigation is framework-agnostic. \\
\midrule

\multicolumn{2}{@{}l}{\textbf{Reviewer \#10}}\\
\midrule
Work reads as a documentation report; add research-methodology / scientific
aspects. &
Added an explicit threat model (\S1), a controlled experimental protocol
with defined ground truth (\S3), quantitative detection metrics (\S4.3),
and a ``Threats to Validity'' subsection (\S6.4). \\

Add a flowchart explaining the method. &
Added Fig.~2, the scanner workflow flowchart. \\

Literature review is short. &
Expanded \S2.2--2.4 and added five peer-reviewed references. \\

References 1, 3, 5--12 are web pages / reports, not authenticated. &
Corrected. Added five verified peer-reviewed references (ECSA 2022,
NDSS 2019 $\times$2, ACM CCS 2013, and a 2025 static-analysis paper), and
removed an unverifiable citation. Remaining web references are limited to
primary sources (framework documentation, OWASP, MITRE CWE, MDN). \\
\midrule

\multicolumn{2}{@{}l}{\textbf{Reviewer \#11}}\\
\midrule
1. Clarify novelty vs.\ Excessive Data Exposure / CWE-200. &
Added \S2.4 ``Novelty and Positioning'' stating three SSR-hydration-specific
claims that distinguish the work from generic CWE-200 and REST-oriented
excessive-exposure analyses. \\

2. Experimental evaluation uses a small synthetic dataset. &
Added the cross-framework fixture corpus and the scalability benchmark;
the controlled synthetic corpus (with exact ground truth) is justified in
\S3 and \S6.4. \\

3. Report precision, recall, $F_1$, and false-positive rate. &
Added the confusion matrix in \S4.3 / Table~6. \\

4. Compare with existing security tools. &
Added \S2.5 and Table~1. \\

5. More peer-reviewed references. &
Added five (see Reviewer \#10). \\

6. Support the multi-framework claim with experiments or state as future
work. &
Supported with experiments: \S4.4 / Table~7 demonstrates detection on
Nuxt, Remix, and SvelteKit. \\
\midrule

\multicolumn{2}{@{}l}{\textbf{Editorial notes}}\\
\midrule
The sources cited in the text were not in the correct order. &
Corrected. References are now numbered in order of first citation, so the
in-text citation numbers ascend sequentially ([1], [2], [3], \dots). \\

English proofreading. &
The manuscript was proofread and standardized to consistent
(American) English throughout. \\

Springer format; author email and affiliation. &
The paper uses the official Springer \texttt{llncs} template; the author's
full affiliation and email are included on the title page. \\

\end{longtable}

\end{document}
